% !TEX root = ../../main-jos.tex
\section{系统总体设计}

\subsection{架构概述}
本文框架采用\textbf{三层协同定向采集架构}，由两个核心角色组成：\textbf{日志节点采集器（Agent）}和\textbf{日志中心（Center）}。Agent 有两种部署形态——网关节点实例和微服务节点实例。网关实例通过 OpenResty Lua 采集 HTTP 流量日志，微服务实例负责采集应用访问日志。Center 的职责包括：接收 Agent 上传的日志、生成关注清单、下发规则、执行二次过滤，最终将结果写入 Loki。总体架构见图~\ref{fig:arch}。

从能力边界来看，整个框架可以分为三层：\textbf{网关预筛选、节点定向采集、中心二次过滤}，如表~\ref{tab:framework_layers} 所示。三层分别负责运行时流量感知、本地日志减量和中心入库控制，每层的输入输出都可以独立审计，也方便与 AIOps 平台的数据采集能力对接。

\begin{table*}[htbp]
  \centering
  \caption{分布式定向日志采集框架的三层输入/输出}
  \label{tab:framework_layers}
  \small
  \resizebox{\textwidth}{!}{%
  \begin{tabular}{p{2.2cm}p{3.0cm}p{3.0cm}p{4.5cm}}
    \toprule
    层级 & 输入 & 输出 & 作用 \\
    \midrule
    网关预筛选 & HTTP access log、状态码、响应时延 & 高价值流量日志与 URL 模式候选 & 从入口流量识别错误、慢请求和热点异常模式 \\
    节点定向采集 & 关注清单、本地 Apache Combined 风格应用日志 & 命中关注清单的应用日志批次 & 在 Sidecar 内完成采集前端减量，降低上传量 \\
    中心二次过滤 & 节点上传批次、关注清单版本、去重状态 & Loki 入库日志与可查询标签 & 统一策略语义，过滤过期或重复日志，保留高价值上下文 \\
    \bottomrule
  \end{tabular}%
  }
\end{table*}

\begin{figure}[htbp]
  \centering
  \includegraphics[width=0.95\textwidth]{fig1_system_overview.pdf}
  \caption{分布式定向日志采集系统总体架构}
  \label{fig:arch}
\end{figure}

\subsection{工作流程}
系统的运行过程分为以下五个步骤：
\begin{enumerate}
  \item 网关在 \texttt{log\_by\_lua} 阶段采集 URL、状态码、响应时延等字段，写入共享内存缓冲区；
  \item 网关 Agent 定期从缓冲区拉取流量日志，通过 gRPC 上传到 Center，Center 将其缓存到 Redis；
  \item Center 对流量日志做高价值过滤和 URL 聚类，生成 Top-$K$ 关注清单，再通过 gRPC 下发给各节点；
  \item 微服务 Agent 拿到清单后，匹配本地应用日志，命中的日志写入固定缓存块，异步压缩上传；
  \item Center 做二次过滤和去重，最后批量推送到 Loki，供 Grafana 查询使用。
\end{enumerate}

\subsection{部署模式}
\textbf{Sidecar 模式}：Agent 与业务容器共享网络命名空间（\texttt{network\_mode: service:*}），适用于 Kubernetes 或 Docker Compose 环境。\textbf{混合开发模式}：基础设施跑在容器里，Center 和 Agent 在本地调试。原型部署在 WSL2 上，网关映射端口 8088（避开 Windows IIS 占用的 80 端口）。
